\documentclass[12pt]{article}

\usepackage[a4paper,margin=1in]{geometry}
\usepackage{amsmath, amssymb, booktabs}
\usepackage{setspace}
\usepackage{enumitem}

\setstretch{1.2}
\setlength{\parindent}{0pt}

\title{\textbf{Final Exam Review Solutions}\\
ECON 204/INTS 211: Contemporary Economic Issues}
\date{}

\begin{document}

\maketitle

%%%%%%%%%%%%%%%%%%%%%%%%%%%%%%%%%%%%%%%%%%%%%%%%%%%%%%
\section*{Part I: Income Distribution and Inequality}

\textbf{Question 1}

Consider the following income distribution:
\[
5{,}000,\; 15{,}000,\; 25{,}000,\; 55{,}000,\; 100{,}000
\]

\subsection*{(a) Total income and mean income}

\textbf{Step 1: Compute total income}
\[
5{,}000 + 15{,}000 + 25{,}000 + 55{,}000 + 100{,}000 = 200{,}000
\]

So total income is:
\[
\boxed{\$200{,}000}
\]

\textbf{Step 2: Compute mean income}

There are 5 households, so:
\[
\bar{y} = \frac{200{,}000}{5} = 40{,}000
\]

Thus, mean income is:
\[
\boxed{\$40{,}000}
\]

\subsection*{(b) Median income}

Since the data are already ordered from smallest to largest, the median is the middle observation.

With 5 values, the middle one is the 3rd income:
\[
\boxed{\$25{,}000}
\]

\subsection*{(c) Income share of the top 20\% and bottom 40\%}

\textbf{Top 20\%}

Since there are 5 households, the top 20\% corresponds to the richest 1 household.

That household earns:
\[
100{,}000
\]

Its income share is:
\[
\frac{100{,}000}{200{,}000} \times 100 = 50\%
\]

So the top 20\% income share is:
\[
\boxed{50\%}
\]

\textbf{Bottom 40\%}

The bottom 40\% corresponds to the poorest 2 households.

Their total income is:
\[
5{,}000 + 15{,}000 = 20{,}000
\]

Their share is:
\[
\frac{20{,}000}{200{,}000} \times 100 = 10\%
\]

So the bottom 40\% income share is:
\[
\boxed{10\%}
\]

\subsection*{(d) Suppose the highest income increases to \$140,000. Recompute the mean and comment on inequality.}

The new income distribution is:
\[
5{,}000,\; 15{,}000,\; 25{,}000,\; 55{,}000,\; 140{,}000
\]

\textbf{Step 1: Compute new total income}
\[
5{,}000 + 15{,}000 + 25{,}000 + 55{,}000 + 140{,}000 = 240{,}000
\]

\textbf{Step 2: Compute new mean}
\[
\bar{y} = \frac{240{,}000}{5} = 48{,}000
\]

So the new mean is:
\[
\boxed{\$48{,}000}
\]

\textbf{Comment on inequality}

The increase in the top income raises the mean from \$40,000 to \$48,000, even though the incomes of the other four households did not change. This indicates that the distribution has become more unequal. The rise in the mean does not reflect a general improvement for all households; instead, it reflects gains concentrated at the top.

\subsection*{(e) Why may the mean be misleading in skewed distributions?}

The mean uses all observations and is heavily affected by very high incomes. In a skewed distribution, a small number of rich households can pull the mean upward. As a result, the mean may suggest that the ``typical'' household is better off than it actually is. The median is often more informative in such cases because it is less sensitive to extreme values.

%%%%%%%%%%%%%%%%%%%%%%%%%%%%%%%%%%%%%%%%%%%%%%%%%%%%%%
\textbf{Question 2 (Gini Coefficient)}

Using the original income distribution:
\[
5{,}000,\; 15{,}000,\; 25{,}000,\; 55{,}000,\; 100{,}000
\]

\subsection*{(a) Construct cumulative population shares and cumulative income shares}

\textbf{Step 1: Total income}

From Question 1:
\[
Y = 200{,}000
\]

\textbf{Step 2: Individual income shares}

\[
\frac{5{,}000}{200{,}000} = 0.025
\]
\[
\frac{15{,}000}{200{,}000} = 0.075
\]
\[
\frac{25{,}000}{200{,}000} = 0.125
\]
\[
\frac{55{,}000}{200{,}000} = 0.275
\]
\[
\frac{100{,}000}{200{,}000} = 0.500
\]

\textbf{Step 3: Cumulative population shares}

There are 5 households, so each household represents 20\% of the population.

\[
X_0 = 0,\quad X_1 = 0.2,\quad X_2 = 0.4,\quad X_3 = 0.6,\quad X_4 = 0.8,\quad X_5 = 1.0
\]

\textbf{Step 4: Cumulative income shares}

\[
Y_0 = 0
\]
\[
Y_1 = 0.025
\]
\[
Y_2 = 0.025 + 0.075 = 0.100
\]
\[
Y_3 = 0.100 + 0.125 = 0.225
\]
\[
Y_4 = 0.225 + 0.275 = 0.500
\]
\[
Y_5 = 0.500 + 0.500 = 1.000
\]

So the cumulative shares are:

\begin{center}
\begin{tabular}{ccc}
\toprule
Household Group & Cumulative Population Share $X_i$ & Cumulative Income Share $Y_i$ \\
\midrule
0 & 0.0 & 0.000 \\
1 & 0.2 & 0.025 \\
2 & 0.4 & 0.100 \\
3 & 0.6 & 0.225 \\
4 & 0.8 & 0.500 \\
5 & 1.0 & 1.000 \\
\bottomrule
\end{tabular}
\end{center}

\subsection*{(b) Compute the Gini coefficient}

Use:
\[
G = 1 - \sum (Y_i + Y_{i-1})(X_i - X_{i-1})
\]

Since the population shares are equally spaced:
\[
X_i - X_{i-1} = 0.2
\]

Now compute each term:

\[
(Y_1 + Y_0)(X_1 - X_0) = (0.025 + 0)(0.2) = 0.005
\]

\[
(Y_2 + Y_1)(X_2 - X_1) = (0.100 + 0.025)(0.2) = 0.125(0.2)=0.025
\]

\[
(Y_3 + Y_2)(X_3 - X_2) = (0.225 + 0.100)(0.2)=0.325(0.2)=0.065
\]

\[
(Y_4 + Y_3)(X_4 - X_3) = (0.500 + 0.225)(0.2)=0.725(0.2)=0.145
\]

\[
(Y_5 + Y_4)(X_5 - X_4) = (1.000 + 0.500)(0.2)=1.500(0.2)=0.300
\]

Now sum:
\[
0.005 + 0.025 + 0.065 + 0.145 + 0.300 = 0.540
\]

Thus:
\[
G = 1 - 0.540 = 0.460
\]

So the Gini coefficient is:
\[
\boxed{0.46}
\]

\subsection*{(c) Interpret the magnitude of the Gini coefficient}

A Gini coefficient of 0 means perfect equality, while a value closer to 1 implies extreme inequality. A Gini of 0.46 suggests a fairly unequal distribution of income. In this small economy, one household earns half of total income, which is consistent with a relatively high level of inequality.

\subsection*{(d) Effect of a transfer from the richest to the poorest household}

If income is transferred from the richest household to the poorest household, the Lorenz curve shifts closer to the 45-degree equality line. This means that cumulative income shares at the bottom of the distribution increase. Since the area between the Lorenz curve and the equality line becomes smaller, the Gini coefficient decreases. Therefore, inequality falls.

%%%%%%%%%%%%%%%%%%%%%%%%%%%%%%%%%%%%%%%%%%%%%%%%%%%%%%
\section*{Part II: Inflation and Cost of Living}

\textbf{Question 3}

Given:
\[
CPI_{2024} = 132.5,\qquad CPI_{2025} = 140.8
\]

Nominal wages rise from:
\[
\$52{,}000 \text{ to } \$54{,}600
\]

\subsection*{(a) Compute the inflation rate}

Use:
\[
\pi = \frac{CPI_{2025} - CPI_{2024}}{CPI_{2024}} \times 100
\]

Substitute:
\[
\pi = \frac{140.8 - 132.5}{132.5} \times 100
\]

Compute the numerator:
\[
140.8 - 132.5 = 8.3
\]

Now divide:
\[
\frac{8.3}{132.5} \approx 0.06264
\]

Convert to percent:
\[
0.06264 \times 100 \approx 6.26\%
\]

So inflation is:
\[
\boxed{6.26\%}
\]

\subsection*{(b) Compute nominal wage growth}

Use:
\[
\text{Nominal wage growth} = \frac{54{,}600 - 52{,}000}{52{,}000} \times 100
\]

\[
54{,}600 - 52{,}000 = 2{,}600
\]

\[
\frac{2{,}600}{52{,}000} = 0.05
\]

\[
0.05 \times 100 = 5\%
\]

So nominal wage growth is:
\[
\boxed{5\%}
\]

\subsection*{(c) Compute real wage growth}

Approximate real wage growth as:
\[
\Delta w_{real} \approx \Delta w_{nominal} - \pi
\]

\[
\Delta w_{real} \approx 5\% - 6.26\% = -1.26\%
\]

So real wage growth is:
\[
\boxed{-1.26\%}
\]

\subsection*{(d) Did purchasing power increase or decrease?}

Because real wage growth is negative, purchasing power fell. Wages rose in nominal terms, but prices rose by more. Therefore, the worker can buy fewer goods and services than before.

\subsection*{(e) How can inflation redistribute income between borrowers and lenders?}

Unexpected inflation benefits borrowers because they repay debt in dollars that are worth less in real terms. Lenders lose because the real value of repayments is lower than expected. Thus, inflation redistributes purchasing power from lenders to borrowers.

%%%%%%%%%%%%%%%%%%%%%%%%%%%%%%%%%%%%%%%%%%%%%%%%%%%%%%
\textbf{Question 4 (Household Inflation)}

\begin{center}
\begin{tabular}{lcc}
\toprule
Category & Weight & Inflation \\
\midrule
Housing & 45\% & 8\% \\
Food & 20\% & 6\% \\
Transport & 15\% & 4\% \\
Other & 20\% & 2\% \\
\bottomrule
\end{tabular}
\end{center}

\subsection*{(a) Compute household-specific inflation}

Use the weighted average formula:
\[
\pi^{HH} = \sum w_i \pi_i
\]

Convert weights to decimals:
\[
0.45,\; 0.20,\; 0.15,\; 0.20
\]

Compute each contribution:

Housing:
\[
0.45 \times 8 = 3.60
\]

Food:
\[
0.20 \times 6 = 1.20
\]

Transport:
\[
0.15 \times 4 = 0.60
\]

Other:
\[
0.20 \times 2 = 0.40
\]

Add them:
\[
3.60 + 1.20 + 0.60 + 0.40 = 5.80
\]

So household-specific inflation is:
\[
\boxed{5.8\%}
\]

\subsection*{(b) Why does this differ from official CPI inflation of 3.5\%?}

This household spends a much larger share on housing, which is the category with the highest inflation. Official CPI is based on an average consumer basket, but individual households have different spending patterns. Therefore, a household with higher housing exposure can face inflation above the national average.

\subsection*{(c) Which households are most affected by this inflation pattern?}

Households that spend large portions of their income on housing and food are most affected. These often include renters, low-income households, single-parent households, and households in expensive urban centres.

%%%%%%%%%%%%%%%%%%%%%%%%%%%%%%%%%%%%%%%%%%%%%%%%%%%%%%
\section*{Part III: Monetary Policy and Interest Rates}

\textbf{Question 5}

Given:
\[
i = 6.5\%, \qquad \pi^e = 3.2\%
\]

\subsection*{(a) Compute the real interest rate}

Using the Fisher approximation:
\[
r \approx i - \pi^e
\]

\[
r \approx 6.5\% - 3.2\% = 3.3\%
\]

So the real interest rate is:
\[
\boxed{3.3\%}
\]

\subsection*{(b) Effect on consumption, investment, and exchange rates}

\textbf{Consumption:}  
A higher real interest rate raises the cost of borrowing. Households reduce spending on credit-financed goods such as cars, appliances, and housing.

\textbf{Investment:}  
Firms face a higher cost of financing projects. Some marginal investment projects become unprofitable, so investment falls.

\textbf{Exchange rates:}  
Higher interest rates can attract foreign capital, increasing demand for the domestic currency. This tends to appreciate the exchange rate, making imports cheaper and exports relatively more expensive.

\subsection*{(c) Effect on inflation expectations}

When the central bank raises interest rates credibly, households and firms may expect lower future inflation. This can reduce wage demands and price-setting behaviour, helping to control current inflation. In this way, monetary policy works partly through expectations.

%%%%%%%%%%%%%%%%%%%%%%%%%%%%%%%%%%%%%%%%%%%%%%%%%%%%%%
\section*{Part IV: Labour Markets and Inequality}

\textbf{Question 6}

\begin{center}
\begin{tabular}{lcc}
\toprule
Group & 2005 & 2025 \\
\midrule
High-skill & \$28 & \$70 \\
Low-skill & \$20 & \$25 \\
\bottomrule
\end{tabular}
\end{center}

\subsection*{(a) Compute the wage gap in both years}

In 2005:
\[
28 - 20 = 8
\]

So the wage gap in 2005 is:
\[
\boxed{\$8}
\]

In 2025:
\[
70 - 25 = 45
\]

So the wage gap in 2025 is:
\[
\boxed{\$45}
\]

\subsection*{(b) Compute percentage wage growth for each group}

\textbf{High-skill workers}

\[
\frac{70 - 28}{28} \times 100
=
\frac{42}{28} \times 100
= 1.5 \times 100 = 150\%
\]

So high-skill wage growth is:
\[
\boxed{150\%}
\]

\textbf{Low-skill workers}

\[
\frac{25 - 20}{20} \times 100
=
\frac{5}{20} \times 100
= 25\%
\]

So low-skill wage growth is:
\[
\boxed{25\%}
\]

\subsection*{(c) Explain using skill-biased technological change}

Skill-biased technological change means that new technologies raise the productivity of highly skilled workers more than that of low-skilled workers. Skilled workers can use technology to become more productive, which increases their wages. At the same time, some low-skill routine tasks are replaced by machines or software, which weakens wage growth for low-skilled workers.

\subsection*{(d) One policy that could reduce inequality}

A strong policy response is investment in education and retraining. If workers gain skills that complement new technologies, they can move into better-paying jobs. Other valid answers include stronger labour protections, earned income tax credits, or improved access to post-secondary education.

%%%%%%%%%%%%%%%%%%%%%%%%%%%%%%%%%%%%%%%%%%%%%%%%%%%%%%
\section*{Part V: Housing Affordability}

\textbf{Question 7}

Income:
\[
\$95{,}000
\]
House price:
\[
\$855{,}000
\]

\subsection*{(a) Compute the price-to-income ratio}

\[
PTI = \frac{855{,}000}{95{,}000} = 9
\]

So:
\[
\boxed{PTI = 9}
\]

\subsection*{(b) Interpret the result}

A PTI of 9 means the house price is nine times annual income. This is generally considered highly unaffordable. Historically, ratios around 3 or 4 were more common. A ratio of 9 suggests severe affordability pressure.

\subsection*{(c) If interest rates increase, how does affordability change even if prices stay constant?}

Higher interest rates increase the monthly mortgage payment for a given loan size. Even if the house price remains unchanged, the cost of financing that house rises. Therefore, affordability worsens because households need a larger monthly budget to carry the mortgage.

%%%%%%%%%%%%%%%%%%%%%%%%%%%%%%%%%%%%%%%%%%%%%%%%%%%%%%
\textbf{Question 8 (Mortgage Burden)}

Mortgage:
\[
\$600{,}000
\]

\subsection*{(a) How do monthly payments change when interest rates rise?}

When interest rates rise, the interest component of the mortgage payment increases. Therefore, monthly mortgage payments increase. The larger the loan, the larger the increase in the monthly payment.

\subsection*{(b) Define and explain a mortgage renewal shock}

A mortgage renewal shock occurs when a household renews its mortgage at a much higher interest rate than before. Even if the outstanding loan has fallen somewhat, the jump in the interest rate can cause a large increase in monthly payments. This creates financial stress and worsens affordability.

\subsection*{(c) Why is this especially important in Canada?}

In Canada, many mortgages have terms of around five years, even though amortization periods are much longer. This means households must renew at prevailing market rates. Therefore, changes in interest rates eventually affect many existing borrowers, making the renewal shock particularly important in Canada.

%%%%%%%%%%%%%%%%%%%%%%%%%%%%%%%%%%%%%%%%%%%%%%%%%%%%%%
\textbf{Question 9 (Rental Market)}

Monthly income:
\[
\$4{,}200
\]

Rent rises from:
\[
\$1{,}300 \text{ to } \$2{,}050
\]

\subsection*{(a) Compute rent-to-income ratios before and after}

\textbf{Before}

\[
\frac{1{,}300}{4{,}200} \times 100 \approx 30.95\%
\]

So the initial ratio is:
\[
\boxed{30.95\%}
\]

\textbf{After}

\[
\frac{2{,}050}{4{,}200} \times 100 \approx 48.81\%
\]

So the new ratio is:
\[
\boxed{48.81\%}
\]

\subsection*{(b) Explain affordability thresholds}

A common rule is that housing becomes unaffordable when shelter costs exceed 30\% of income. Some analysts use 50\% to indicate severe housing stress. In this example, the renter moves from just above the 30\% threshold to nearly 49\%, indicating very serious affordability pressure.

\subsection*{(c) How do rising rents contribute to inequality?}

Rising rents reduce the ability of renters to save and accumulate wealth. Homeowners may benefit from rising home values, while renters face higher housing costs without gaining assets. This widens inequality between renters and owners, and often between younger and older households.

%%%%%%%%%%%%%%%%%%%%%%%%%%%%%%%%%%%%%%%%%%%%%%%%%%%%%%
\section*{Part VI: Integrated Policy Analysis}

\textbf{Question 10}

Policy: zoning reform to increase housing supply.

\subsection*{(a) Short-run and long-run effects on housing prices}

\textbf{Short run:}  
In the short run, prices may not change much because construction takes time and supply cannot expand immediately.

\textbf{Long run:}  
In the long run, zoning reform can allow more housing units to be built. As supply expands, price growth should slow, and affordability may improve.

\subsection*{(b) Effect on inequality}

Increasing housing supply can reduce exclusion from high-demand areas and help middle-income households access housing. If more rental housing is built, rent pressures may ease, benefiting lower-income households. Thus, zoning reform can reduce some dimensions of housing inequality.

\subsection*{(c) One unintended consequence}

One unintended consequence could be neighbourhood opposition or infrastructure strain. If roads, schools, and transit are not expanded, denser development may create congestion or reduce service quality.

%%%%%%%%%%%%%%%%%%%%%%%%%%%%%%%%%%%%%%%%%%%%%%%%%%%%%%
\textbf{Question 11 (Integrated Analysis)}

Suppose inflation rises to 6\% and the central bank raises interest rates significantly.

\subsection*{(a) Impact on housing demand, mortgage payments, and rental markets}

\textbf{Housing demand:}  
Higher interest rates reduce borrowing capacity, so fewer households can afford to buy homes. Housing demand falls.

\textbf{Mortgage payments:}  
Existing borrowers renewing at higher rates face higher monthly payments. Variable-rate borrowers may also see payments rise more quickly.

\textbf{Rental markets:}  
As some potential buyers are priced out of ownership, they remain renters or enter the rental market. This raises rental demand and may push rents higher.

\subsection*{(b) Why might inflation fall but housing affordability worsen?}

Higher interest rates reduce aggregate demand and help control inflation. However, they also increase borrowing costs and mortgage payments. Therefore, even if inflation begins to decline, housing affordability can worsen because financing housing becomes more expensive.

\subsection*{(c) Distributional effects across income groups}

Low- and middle-income households are usually hit hardest. They spend larger shares of income on housing and have less financial cushion. Existing homeowners with small or no mortgages may be less affected than new buyers and renters. Younger households often face the greatest burden because they are more likely to rent or hold large recent mortgages.

%%%%%%%%%%%%%%%%%%%%%%%%%%%%%%%%%%%%%%%%%%%%%%%%%%%%%%
\section*{Concluding Note}

These review problems combine numerical calculations, interpretation, and policy analysis. On an exam, it is important not only to calculate correctly, but also to explain clearly what the numbers mean economically. Strong answers connect the quantitative result to the broader issues of inequality, inflation, labour markets, housing, and policy trade-offs.

\end{document}